\documentclass[a4paper,9pt]{extarticle}

%=========================================================
% PACKAGES
%=========================================================

\usepackage[
a4paper,
top=1.8cm,
bottom=1.5cm,
footskip=7mm,
left=1cm,
right=1cm,
headheight=14pt,
headsep=5mm
]{geometry}
\usepackage[T1]{fontenc}
\usepackage[utf8]{inputenc}
\usepackage[french]{babel}

\usepackage[sfdefault]{FiraSans}
\renewcommand{\familydefault}{\sfdefault}

\usepackage{multicol}
\setlength{\columnsep}{0.7cm}

\usepackage{amsmath,amssymb}

\usepackage{titlesec}
\usepackage{enumitem}

\usepackage{tcolorbox}
\tcbuselibrary{skins,breakable}

\usepackage{xcolor}

\usepackage{graphicx}
\usepackage{float}

\usepackage{tikz}

\usepackage{fancyhdr}
\usepackage[hidelinks]{hyperref}


%=========================================================
% PAGE SETTINGS
%=========================================================

\setlength{\parindent}{0pt}
\setlength{\parskip}{1.5pt}

\setlength{\abovedisplayskip}{4pt}
\setlength{\belowdisplayskip}{4pt}

\setlist[itemize]{
leftmargin=*,
itemsep=1pt,
topsep=1pt
}


%=========================================================
% COLORS
%=========================================================

\definecolor{Blue}{HTML}{E3F0FF}
\definecolor{Green}{HTML}{E3F7EA}
\definecolor{Pink}{HTML}{FFE5ED}
\definecolor{Orange}{HTML}{FFF0DC}
\definecolor{Purple}{HTML}{EEE7FF}
\definecolor{Yellow}{HTML}{FFFADB}


\definecolor{BlueTitle}{HTML}{285B8F}
\definecolor{GreenTitle}{HTML}{287052}
\definecolor{PinkTitle}{HTML}{A84063}
\definecolor{OrangeTitle}{HTML}{A96500}
\definecolor{PurpleTitle}{HTML}{6245A5}


%=========================================================
% HEADER
%=========================================================

\pagestyle{fancy}

\fancyhf{}

\fancyhead[L]{\small\textbf{Computer Security}}
\fancyhead[C]{\small Introduction}
\fancyhead[R]{\small Week 01 --- 07-13/09/2026}

\fancyfoot[C]{\small\href{https://github.com/AlphaXT38G/notes_ba5}{github.com/AlphaXT38G/notes\_ba5}}
\fancyfoot[R]{\small \thepage}

\renewcommand{\headrulewidth}{0.4pt}



%=========================================================
% TITLES
%=========================================================

\titleformat{\section}
{\large\bfseries\color{BlueTitle}}
{}{0pt}{}

\titlespacing{\section}
{0pt}{6pt}{3pt}


% separator

\newcommand{\separator}{
\vspace{2pt}
\hrule
\vspace{3pt}
}


%=========================================================
% BOXES
%=========================================================

% DEFINITION - blue
\newtcolorbox{definition}{
enhanced,
breakable,
colback=Blue,
colframe=BlueTitle,
boxrule=0.7pt,
arc=0mm,
title=\textbf{Definition},
fonttitle=\bfseries
}

% IMPORTANT - pink
\newtcolorbox{important}{
enhanced,
breakable,
colback=Pink,
colframe=PinkTitle,
boxrule=0.7pt,
arc=0mm,
title=\textbf{Key idea},
fonttitle=\bfseries
}

% FORMULAS - green
\newtcolorbox{formula}{
enhanced,
breakable,
colback=Green,
colframe=GreenTitle,
boxrule=0.7pt,
arc=0mm,
title=\textbf{Key formulas},
fonttitle=\bfseries
}

% WARNING - orange
\newtcolorbox{warning}{
enhanced,
breakable,
colback=Orange,
colframe=OrangeTitle,
boxrule=0.7pt,
arc=0mm,
title=\textbf{Warning},
fonttitle=\bfseries
}

% EXAMPLE - purple
\newtcolorbox{example}{
enhanced,
breakable,
colback=Purple,
colframe=PurpleTitle,
boxrule=0.7pt,
arc=0mm,
title=\textbf{Example},
fonttitle=\bfseries
}

% SUMMARY - orange
\newtcolorbox{summary}{
enhanced,
breakable,
colback=Yellow,
colframe=OrangeTitle,
boxrule=0.7pt,
arc=0mm,
title=\textbf{Summary},
fonttitle=\bfseries
}



%=========================================================
% FIGURE BOX
%=========================================================

\newtcolorbox{figurebox}{
enhanced,
colback=white,
colframe=gray,
boxrule=0.5pt,
arc=0mm,
center
}



%=========================================================
% DOCUMENT
%=========================================================

\begin{document}

%\footnotesize   --> To reduce the size of the text

\thispagestyle{fancy}



\begin{multicols}{2}


%=========================================================
\section{Computer Security}


\begin{definition}

\textbf{COMPUTER SECURITY}

Properties of a computer system must hold
in the presence of a resourced strategic adversary.

\end{definition}


\begin{important}

If the properties hold within the threat model, it means that there are no vulnerabilities that can be exploited to materialize threats, and no harm happens.

\end{important}



%=========================================================
\section{Properties}


\textbf{What Properties?}


\textbf{TRADITIONAL PROPERTIES}

\begin{itemize}

\item \textbf{Confidentiality} -- prevention of unauthorized disclosure of information

(e.g. The adversary should not be able to read my bank statement)

\item \textbf{Integrity} -- prevention of unauthorized modification of information

(e.g. The adversary should not be able to change my bank balance)

\item \textbf{Availability} -- prevention of unauthorized denial of service

(e.g. The adversary should not prevent me accessing my bank account)

\end{itemize}


\separator


\textbf{NOT-SO-TRADITIONAL (BUT IMPORTANT) PROPERTIES}

\begin{itemize}

\item Authenticity

\item Anonymity

\item Isolation

\item Non-repudiation

\item \dots

\end{itemize}



%=========================================================
\section{Security Policy}


\begin{definition}

\textbf{SECURITY POLICY:} a high-level description of the security properties that must
hold in the system in relation to assets and principals.

\end{definition}


\begin{definition}

\textbf{Assets (objects):} anything with value (e.g., data, files, memory) that needs to be protected.

\end{definition}


\begin{definition}

\textbf{Principals (subjects):} people, computer programs, services,\dots{} (may not contain the adversary)

\end{definition}



%=========================================================
\section{Threat Model and Adversary}


\begin{definition}

\textbf{THREAT MODEL:} describes the resources available to the adversary and the adversary's capabilities (observe, influence, corrupt,\dots)

\end{definition}


\begin{definition}

The \textbf{adversary} is a malicious entity aiming at breaching the security policy. The adversary is strategic: the adversary will choose the optimal way to use her resources to mount an attack that violates the security properties.

\end{definition}


\separator


\begin{definition}

\textbf{THREAT MODEL}

Technical term to define the
adversary's capabilities.

\end{definition}


\begin{example}

\begin{itemize}

\item The adversary can observe my connection

\item The adversary can corrupt my machine

\item The adversary controls a bank employee

\end{itemize}

\end{example}


\begin{definition}

\textbf{VULNERABILITY}

Specific weakness that could be exploited by
adversaries with interest in a lot of different
assets

\end{definition}


\begin{example}

\begin{itemize}

\item The banking API is not protected

\item The password appears in plain text in my screen

\end{itemize}

\end{example}


\begin{definition}

\textbf{THREAT}

What is the feared event, the goal of the
adversary that we don't want
materialized

\end{definition}


\begin{example}

\begin{itemize}

\item A hacker wants to retrieve money breaking
into the bank's system

\item A student wants to learn my password by
looking over my shoulder

\end{itemize}

\end{example}


\begin{definition}

\textbf{HARM}

The bad thing that happens when the
threat materializes

\end{definition}


\begin{example}

\begin{itemize}

\item The adversary steals money

\item The adversary blocks access to the bank

\item The adversary learns my password

\item The adversary reads the message

\end{itemize}

\end{example}



%=========================================================
\section{Security Mechanism}


\begin{definition}

\textbf{SECURITY MECHANISM:} Technical mechanism used to ensure that the security policy
is not violated by an adversary within the threat model.

\end{definition}


\begin{example}

\textbf{Example 1}

\textbf{Policy:} ensure the log of transactions is not tampered with by a single employee

\textbf{Mechanism:} keep a copy of the log on multiple computers, such that no single employee has access to all of them

\end{example}



%=========================================================
\section{Attacker vs Defender}


\begin{definition}

\textbf{ATTACKER} Just one way to violate one security property is enough!
(within the threat model)

\end{definition}


\begin{definition}

\textbf{DEFENDER} No adversary strategy can violate the security policy

\end{definition}


\begin{important}

A system is ``secure'' if an adversary constrained by a
specific threat model cannot violate the security policy

\end{important}


\separator


\begin{definition}

\textbf{SECURITY ARGUMENT:} rigorous argument that the security mechanisms in place are
indeed effective in maintaining the security policy (verbal or mathematical).

Subject to the assumptions of the threat model.

\end{definition}


\begin{warning}

\textbf{FOR A THREAT MODEL TO BE USEFUL} The model must constrain the adversary, otherwise we cannot
make a security argument

\end{warning}



%=========================================================
\section{Steps of Security Engineering}


\begin{summary}

\textbf{1.- High-level specification}

\begin{itemize}

\item Define the architecture of the system (e.g., high level block diagram)

\item Define the security policy (principals, assets, security properties)

\item Define the threat model

\end{itemize}

\textbf{2.- Security design}

\begin{itemize}

\item Select / Design security mechanisms

\item State your security argument: which controls maintain which
properties?

\end{itemize}

\textbf{3.- Secure implementation}

\begin{itemize}

\item Implement mechanisms

\item Ensure they conform to the design model

\item Security testing

\end{itemize}

\end{summary}



\end{multicols}


\end{document}
